\setcounter{chapter}{4}
\chapter{Lebesgue測度}

\paragraph{本章の目標}
Lebesgue外測度 $\lambda^*$ とLebesgue内測度 $\lambda_*$ が一致する集合を
Lebesgue可測集合と定める．
そのうえで，Lebesgue可測集合全体が $\sigma$-加法族をなし，
Lebesgue測度が可算加法的な測度になることを確認する．
最後に，内外正則性とCarathéodory条件による判定法を述べる．

\section{Lebesgue可測集合とLebesgue測度}

\begin{definition}[Lebesgue可測集合とLebesgue測度]
集合 $A \subset \RR^d$ がLebesgue可測であるとは
\[
\lambda_*(A)=\lambda^*(A)
\]
が成り立つことをいう．
Lebesgue可測集合全体を
\[
\calL^d:=\{A\subset\RR^d\mid \lambda_*(A)=\lambda^*(A)\}
\]
と書く．
$A\in\calL^d$ のとき，この共通値を
\[
\lambda(A):=\lambda_*(A)=\lambda^*(A)
\]
と書き，$A$ のLebesgue測度という．
\end{definition}

\begin{definition}[Lebesgue零集合]
集合 $N\subset\RR^d$ が
\[
\lambda^*(N)=0
\]
を満たすとき，$N$ をLebesgue零集合という．
\end{definition}

\begin{theorem}[Lebesgue零集合]
Lebesgue零集合はLebesgue可測であり，Lebesgue測度は $0$ である．
\end{theorem}
\begin{proof}
第4章で示した基本性質より
\[
0\le \lambda_*(N)\le \lambda^*(N)
\]
である．
$\lambda^*(N)=0$ なら
\[
\lambda_*(N)=\lambda^*(N)=0
\]
だから，$N$ はLebesgue可測で $\lambda(N)=0$ である．
\end{proof}

\begin{corollary}[可算集合]
一点集合，可算集合，$\QQ^d$，$\ZZ^d$ はLebesgue可測であり，測度 $0$ をもつ．
特に1次元では $\QQ$ と $\ZZ$ はLebesgue測度 $0$ である．
\end{corollary}
\begin{proof}
第3章で，一点集合および可算集合のLebesgue外測度は $0$ であることを示した．
したがって前定理より従う．
\end{proof}

\begin{theorem}[Jordan可測集合]
Jordan可測集合 $A\subset\RR^d$ はLebesgue可測であり，
\[
\lambda(A)=m_J(A)
\]
が成り立つ．
\end{theorem}
\begin{proof}
第3章で，Jordan可測集合に対して
\[
\lambda^*(A)=m_J(A)
\]
が成り立つことを示した．

$\eps>0$ を取る．
Jordan可測性の近似定理と，基本集合をコンパクト集合で内側から近似できることにより，
あるコンパクト集合 $K\subset A$ が存在して
\[
m_J(K)>m_J(A)-\eps
\]
を満たす．
コンパクト集合はJordan可測なので，第3章の整合性より
\[
\lambda^*(K)=m_J(K)
\]
である．
したがって内測度の定義から
\[
\lambda_*(A)\ge \lambda^*(K)=m_J(K)>m_J(A)-\eps
=\lambda^*(A)-\eps
\]
である．
一方，常に $\lambda_*(A)\le \lambda^*(A)$ だから，
$\eps>0$ の任意性より
\[
\lambda_*(A)=\lambda^*(A)
\]
を得る．
よって $A$ はLebesgue可測であり，
\[
\lambda(A)=\lambda^*(A)=m_J(A)
\]
である．
\end{proof}

\begin{example}
\[
\QQ\cap[0,1]
\]
は可算集合なのでLebesgue可測であり，
\[
\lambda_1(\QQ\cap[0,1])=0
\]
である．
一方，$[0,1]$ はJordan可測で $m_J([0,1])=1$ だから
\[
\lambda_1([0,1])=1
\]
である．
後で示す有限加法性より
\[
\lambda_1([0,1]\setminus\QQ)=1
\]
となる．
\end{example}

\section{Carathéodory条件}

\begin{definition}[Carathéodory条件]
集合 $A\subset\RR^d$ がLebesgue外測度 $\lambda^*$ に関するCarathéodory条件を満たすとは，
任意の $E\subset\RR^d$ に対して
\[
\lambda^*(E)
=
\lambda^*(E\cap A)+\lambda^*(E\setminus A)
\]
が成り立つことをいう．
\end{definition}

\begin{theorem}[作業用の判定法]
集合 $A\subset\RR^d$ について，次は同値である．
\begin{enumerate}
    \item $A$ はLebesgue可測である
    \item $A$ はCarathéodory条件を満たす
\end{enumerate}
\end{theorem}
\begin{proof}
本章末の判定法で，正則性を含む形で証明する．
証明の核は第6章で述べるCarathéodoryの定理である．
\end{proof}

\section{Lebesgue可測集合の性質}

\begin{theorem}[Lebesgue可測集合の $\sigma$-加法族性]
Lebesgue可測集合全体 $\calL^d$ は $\RR^d$ 上の $\sigma$-加法族である．
すなわち，次が成り立つ．
\begin{enumerate}
    \item $\RR^d\in\calL^d$
    \item $A\in\calL^d$ ならば $\RR^d\setminus A\in\calL^d$
    \item $\{A_n\}_{n=1}^{\infty}\subset\calL^d$ ならば
    \[
    \bigcup_{n=1}^{\infty}A_n\in\calL^d
    \]
\end{enumerate}
\end{theorem}
\begin{proof}
作業用の判定法により，$\calL^d$ は
$\lambda^*$ に関するCarathéodory可測集合全体と一致する．
第6章のCarathéodoryの定理によれば，外測度に関するCarathéodory可測集合全体は
$\sigma$-加法族である．
したがって $\calL^d$ も $\sigma$-加法族である．
\end{proof}

\begin{corollary}[集合演算]
$A,B\in\calL^d$ ならば
\[
A\cup B,\qquad A\cap B,\qquad A\setminus B
\]
はいずれもLebesgue可測である．
また $\{A_n\}_{n=1}^{\infty}\subset\calL^d$ ならば
\[
\bigcap_{n=1}^{\infty}A_n
\]
もLebesgue可測である．
\end{corollary}
\begin{proof}
$\sigma$-加法族の基本性質から従う．
\end{proof}

\section{Lebesgue測度の性質}

\begin{lemma}[有限加法性]
$A,B\in\calL^d$ が互いに素ならば
\[
\lambda(A\cup B)=\lambda(A)+\lambda(B)
\]
が成り立つ．
\end{lemma}
\begin{proof}
$A$ はCarathéodory条件を満たす．
$E=A\cup B$ としてその条件を適用すると，
\[
\lambda^*(A\cup B)
=
\lambda^*((A\cup B)\cap A)+\lambda^*((A\cup B)\setminus A)
=
\lambda^*(A)+\lambda^*(B)
\]
である．
$A,B,A\cup B$ はLebesgue可測なので，これは
\[
\lambda(A\cup B)=\lambda(A)+\lambda(B)
\]
を意味する．
\end{proof}

\begin{lemma}[互いに素な列への分解]
Lebesgue可測集合列 $\{A_n\}_{n=1}^{\infty}$ に対して
\[
B_1:=A_1,\qquad
B_n:=A_n\setminus\bigcup_{k=1}^{n-1}A_k\quad(n\ge2)
\]
とおくと，$\{B_n\}$ は互いに素なLebesgue可測集合列であり，
\[
\bigcup_{n=1}^{\infty}A_n=\bigsqcup_{n=1}^{\infty}B_n,
\qquad
B_n\subset A_n
\]
が成り立つ．
\end{lemma}
\begin{proof}
$\calL^d$ は $\sigma$-加法族なので，各 $B_n$ はLebesgue可測である．
定義から互いに素性と表示は明らかである．
\end{proof}

\begin{theorem}[完全加法性]
互いに素なLebesgue可測集合列 $\{A_n\}_{n=1}^{\infty}$ に対して
\[
\lambda\left(\bigsqcup_{n=1}^{\infty}A_n\right)
=
\sum_{n=1}^{\infty}\lambda(A_n)
\]
が成り立つ．
\end{theorem}
\begin{proof}
作業用の判定法により，Lebesgue可測集合はCarathéodory可測集合である．
第6章のCarathéodoryの定理によれば，外測度 $\lambda^*$ を
Carathéodory可測集合全体に制限したものは測度である．
したがって互いに素な可測集合列に対して可算加法性が成り立つ．
\end{proof}

\begin{theorem}[基本性質]
Lebesgue測度について次が成り立つ．
\begin{enumerate}
    \item $\lambda(\emptyset)=0$
    \item 単調性：$A,B\in\calL^d$ かつ $A\subset B$ ならば
    \[
    \lambda(A)\le \lambda(B)
    \]
    \item 完全劣加法性：$\{A_n\}_{n=1}^{\infty}\subset\calL^d$ ならば
    \[
    \lambda\left(\bigcup_{n=1}^{\infty}A_n\right)
    \le
    \sum_{n=1}^{\infty}\lambda(A_n)
    \]
\end{enumerate}
\end{theorem}
\begin{proof}
\begin{enumerate}
    \item $\emptyset$ はLebesgue零集合であるから可測で，測度 $0$ をもつ．

    \item $A\subset B$ とする．
    このとき
    \[
    B=A\sqcup(B\setminus A)
    \]
    であり，有限加法性より
    \[
    \lambda(B)=\lambda(A)+\lambda(B\setminus A)\ge \lambda(A)
    \]
    である．

    \item 前補題の記法を用いると
    \[
    \bigcup_{n=1}^{\infty}A_n=\bigsqcup_{n=1}^{\infty}B_n,
    \qquad
    B_n\subset A_n
    \]
    である．
    完全加法性と単調性より
    \[
    \lambda\left(\bigcup_{n=1}^{\infty}A_n\right)
    =
    \sum_{n=1}^{\infty}\lambda(B_n)
    \le
    \sum_{n=1}^{\infty}\lambda(A_n)
    \]
    を得る．
\end{enumerate}
\end{proof}

\begin{theorem}[下からの連続性]
Lebesgue可測集合の単調増大列
\[
A_1\subset A_2\subset \cdots
\]
に対して
\[
\lambda\left(\bigcup_{n=1}^{\infty}A_n\right)
=
\lim_{n\to\infty}\lambda(A_n)
\]
が成り立つ．
\end{theorem}
\begin{proof}
\[
B_1:=A_1,\qquad B_n:=A_n\setminus A_{n-1}\quad(n\ge2)
\]
とおく．
すると $\{B_n\}$ は互いに素で，
\[
A_N=\bigsqcup_{n=1}^{N}B_n,\qquad
\bigcup_{n=1}^{\infty}A_n=\bigsqcup_{n=1}^{\infty}B_n
\]
である．
完全加法性より
\[
\lambda(A_N)=\sum_{n=1}^{N}\lambda(B_n),
\qquad
\lambda\left(\bigcup_{n=1}^{\infty}A_n\right)
=
\sum_{n=1}^{\infty}\lambda(B_n).
\]
$N\to\infty$ とすれば主張を得る．
\end{proof}

\begin{theorem}[上からの連続性]
Lebesgue可測集合の単調減少列
\[
A_1\supset A_2\supset \cdots
\]
に対して，$\lambda(A_1)<\infty$ ならば
\[
\lambda\left(\bigcap_{n=1}^{\infty}A_n\right)
=
\lim_{n\to\infty}\lambda(A_n)
\]
が成り立つ．
\end{theorem}
\begin{proof}
\[
B_n:=A_1\setminus A_n
\]
とおくと，
\[
B_1\subset B_2\subset \cdots,
\qquad
\bigcup_{n=1}^{\infty}B_n
=
A_1\setminus\bigcap_{n=1}^{\infty}A_n
\]
である．
下からの連続性より
\[
\lambda\left(A_1\setminus\bigcap_{n=1}^{\infty}A_n\right)
=
\lim_{n\to\infty}\lambda(A_1\setminus A_n)
\]
である．
$A_n\subset A_1$ かつ $\lambda(A_1)<\infty$ なので有限加法性より
\[
\lambda(A_1\setminus A_n)=\lambda(A_1)-\lambda(A_n)
\]
である．
同様に
\[
\lambda\left(A_1\setminus\bigcap_{n=1}^{\infty}A_n\right)
=
\lambda(A_1)-\lambda\left(\bigcap_{n=1}^{\infty}A_n\right).
\]
これらを合わせれば主張が従う．
\end{proof}

\begin{proposition}[平行移動不変性]
$A\in\calL^d$ と $c\in\RR^d$ に対して
\[
A+c:=\{x+c\mid x\in A\}
\]
もLebesgue可測であり，
\[
\lambda(A+c)=\lambda(A)
\]
が成り立つ．
\end{proposition}
\begin{proof}
第3章でLebesgue外測度の平行移動不変性を示した．
また $K$ がコンパクトなら $K+c$ もコンパクトなので，内測度も平行移動で不変である．
したがって
\[
\lambda^*(A+c)=\lambda^*(A),
\qquad
\lambda_*(A+c)=\lambda_*(A).
\]
$A$ がLebesgue可測なら $A+c$ もLebesgue可測で，測度は変わらない．
\end{proof}

\section{Lebesgue測度の正則性}

\begin{theorem}[コンパクト集合]
コンパクト集合 $K\subset\RR^d$ はLebesgue可測であり，
\[
\lambda(K)=\lambda^*(K)
\]
が成り立つ．
\end{theorem}
\begin{proof}
第4章で
\[
\lambda_*(K)=\lambda^*(K)
\]
を示した．
したがって $K$ はLebesgue可測で，その測度は共通値 $\lambda^*(K)$ に等しい．
\end{proof}

\begin{lemma}[座標超平面片]
有界閉区間 $I\subset\RR^d$ と座標 $i$，実数 $c$ に対して
\[
F:=\{x\in I\mid x_i=c\}
\]
はLebesgue零集合である．
\end{lemma}
\begin{proof}
$F$ は任意に薄い直方体で覆える．
具体的には，$\eps>0$ に対して幅 $2\delta$ の板で $F$ を覆い，
$\delta>0$ を十分小さく取ればその体積を $\eps$ 未満にできる．
したがって $\lambda^*(F)=0$ である．
\end{proof}

\begin{theorem}[開集合]
開集合 $G\subset\RR^d$ はLebesgue可測である．
\end{theorem}
\begin{proof}
まず有界開区間
\[
J=(a_1,b_1)\times\cdots\times(a_d,b_d)
\]
を考える．
同じ端点をもつ閉区間を $\overline J$ とすると，
$\overline J$ はコンパクトだからLebesgue可測である．
また
\[
\overline J\setminus J
\]
は有限個の座標超平面片の合併に含まれるためLebesgue零集合である．
よって
\[
J=\overline J\setminus(\overline J\setminus J)
\]
はLebesgue可測である．

端点がすべて有理数である有界開区間全体は可算であり，$\RR^d$ の開集合の基底をなす．
したがって任意の開集合 $G$ は，そのような有界開区間の可算和として書ける．
$\calL^d$ は $\sigma$-加法族だから，$G$ はLebesgue可測である．
\end{proof}

\begin{theorem}[内正則性]
任意のLebesgue可測集合 $A\subset\RR^d$ に対して
\[
\lambda(A)
=
\sup\{\lambda(K)\mid K\subset A,\ K\text{ はコンパクト}\}
\]
が成り立つ．
\end{theorem}
\begin{proof}
$A$ がLebesgue可測なら
\[
\lambda(A)=\lambda_*(A)
\]
である．
内測度の定義より
\[
\lambda_*(A)
=
\sup\{\lambda^*(K)\mid K\subset A,\ K\text{ はコンパクト}\}.
\]
コンパクト集合はLebesgue可測なので $\lambda^*(K)=\lambda(K)$ と書ける．
よって主張が従う．
\end{proof}

\begin{theorem}[外正則性]
任意のLebesgue可測集合 $A\subset\RR^d$ に対して
\[
\lambda(A)
=
\inf\{\lambda(G)\mid A\subset G,\ G\text{ は開集合}\}
\]
が成り立つ．
\end{theorem}
\begin{proof}
$A$ がLebesgue可測なら
\[
\lambda(A)=\lambda^*(A)
\]
である．
第3章の開集合による外正則性より
\[
\lambda^*(A)
=
\inf\{\lambda^*(G)\mid A\subset G,\ G\text{ は開集合}\}
\]
である．
開集合はLebesgue可測だから $\lambda^*(G)=\lambda(G)$ と書ける．
したがって主張を得る．
\end{proof}

\section{Lebesgue可測集合の判定法}

\begin{theorem}[判定法]
$A\subset\RR^d$ とする．
次は同値である．
\begin{enumerate}
    \item $A$ はLebesgue可測である
    \item $A$ はCarathéodory条件を満たす
    \item 任意の $R>0$ と $\eps>0$ に対し，開集合 $G$ とコンパクト集合 $K$ が存在して
    \[
    K\subset A\cap[-R,R]^d\subset G,
    \qquad
    \lambda^*(G\setminus K)<\eps
    \]
    を満たす
\end{enumerate}
さらに $\lambda^*(A)<\infty$ のときは，これらは次の条件とも同値である．
\begin{enumerate}
    \item[(4)] 任意の $\eps>0$ に対して開集合 $G$ とコンパクト集合 $K$ が存在して
    \[
    K\subset A\subset G,
    \qquad
    \lambda^*(G\setminus K)<\eps
    \]
    を満たす
\end{enumerate}
\end{theorem}
\begin{proof}
(1) と (2) の同値は，第6章のCarathéodoryの定理をLebesgue外測度 $\lambda^*$ に適用し，
Lebesgue流の可測性 $\lambda_*(A)=\lambda^*(A)$ と照合することで得られる．

(1) から (3) を示す．
\[
A_R:=A\cap[-R,R]^d
\]
とおく．
$A_R$ はLebesgue可測で，しかも有限測度である．
内正則性より，コンパクト集合 $K\subset A_R$ で
\[
\lambda(A_R)-\lambda(K)<\frac{\eps}{2}
\]
を満たすものが取れる．
外正則性より，開集合 $G\supset A_R$ で
\[
\lambda(G)-\lambda(A_R)<\frac{\eps}{2}
\]
を満たすものが取れる．
すると $K\subset G$ であり，
\[
\lambda^*(G\setminus K)
=
\lambda(G\setminus K)
=
\lambda(G)-\lambda(K)
<\eps
\]
である．

(3) から (1) を示す．
各 $R>0$ について $A_R=A\cap[-R,R]^d$ とおく．
(3) により，任意の $\eps>0$ に対して
\[
K\subset A_R\subset G,
\qquad
\lambda^*(G\setminus K)<\eps
\]
を満たすコンパクト集合 $K$ と開集合 $G$ が取れる．
このとき $K$ と $G$ はLebesgue可測であり，
\[
\lambda^*(A_R)-\lambda_*(A_R)
\le
\lambda(G)-\lambda(K)
=
\lambda(G\setminus K)
<\eps
\]
である．
$\eps>0$ は任意だから $A_R$ はLebesgue可測である．
したがって
\[
A=\bigcup_{n=1}^{\infty}\bigl(A\cap[-n,n]^d\bigr)
\]
はLebesgue可測である．

最後に $\lambda^*(A)<\infty$ と仮定する．
(1) から (4) は，内正則性と外正則性を直接 $A$ に適用すればよい．
(4) から (1) は，上の (3) から (1) の証明と同じ議論を $A$ そのものに適用すれば従う．
\end{proof}

\begin{remark}
条件 (4) の形は有限測度の集合に対する判定法である．
例えば $A=\RR^d$ の場合，どのコンパクト集合 $K$ を取っても
\[
\RR^d\setminus K
\]
は無限の測度をもつので，(4) は成り立たない．
一般のLebesgue可測集合に対しては，条件 (3) のように有界な窓で切った局所版を用いる．
\end{remark}
